\chapter{Embedding, Complementation, and Relativisation}\label{ch:embedding}

\section{The conjunctive / quotative linker}\label{sec:conjunctive-linker}

Iridian has a verbal linking form in \ird{-e}, used in clause linking and
quotative constructions. The suffix selects the front grade at the
right edge of the fully inflected verb. Its morphological grade is independent
of the fact that its pronounced vowel is /ɛ/.

\subsection{Forms and morphophonology}

The finite and linked forms exhibit both productive consonant alternations
and construction-specific ending correspondences:

\pex
\a \ird{pjášček} $\rightarrow$ \ird{pjáščice}
\a \ird{pjáščav} $\rightarrow$ \ird{pjáščavě}
\a \ird{mužna} $\rightarrow$ \ird{mužne}
\a \ird{pjaštíme} $\rightarrow$ \ird{pjaštímně}
\a \ird{mužnáž} $\rightarrow$ \ird{mužnáže}
\xe

\begin{table}[ht]
  \footnotesize
  \caption{Representative linker outputs across finite paradigms.}
  \label{tab:linker-paradigm}
  \begin{tabularx}{\textwidth}{l l l l}
    \toprule
    \textsc{mood} & \textsc{domain / category} & \textsc{finite form} & \textsc{linked form} \\
    \midrule
    Indicative & perfective & \ird{pjaštník} & \ird{pjaštníce} \\
    Indicative & retrospective & \ird{pjaštnáv} & \ird{pjaštnávě} \\
    Indicative & continuous & \ird{mužna} & \ird{mužne} \\
    Indicative & progressive & \ird{pjaštíme} & \ird{pjaštímně} \\
    Indicative & contemplative & \ird{mužnáž} & \ird{mužnáže} \\
    Imperative & IPFV & \ird{pjaštím} & \ird{pjaštímě} \\
    Imperative & PFV & \ird{pjaštovím} & \ird{pjaštovímě} \\
    Prohibitive & IPFV & \ird{pjašté} & \ird{pjaště} \\
    Prohibitive & PFV & \ird{pjaštové} & \ird{pjaštově} \\
    Hortative & IPFV & \ird{pjaštká} & \ird{pjaštce} \\
    Hortative & PFV & \ird{pjaštovká} & \ird{pjaštovce} \\
    Negative hortative & IPFV & \ird{pjaštku} & \ird{pjaštcuje} \\
    Negative hortative & PFV & \ird{pjaštovku} & \ird{pjaštovcuje} \\
    Irrealis & IPFV & \ird{pjáščev} & \ird{pjáščevě} \\
    Irrealis & PFV & \ird{pjaštóvěv} & \ird{pjaštóvěvě} \\
    Negative irrealis & IPFV & \ird{pjaštir} & \ird{pjaštire} \\
    Negative irrealis & PFV & \ird{pjaštovir} & \ird{pjaštovire} \\
    \bottomrule
  \end{tabularx}
\end{table}

Consonant-final forms expose the grade directly: /k/ selects /ts/,
input /ts/ selects /tʃ/, and final labials acquire secondary palatalisation.
Before linker /ɛ/, \ird{tě\,dě\,ně\,lě} represent the corresponding
soft-coronal sequences.
The paradigm also contains listed ending allomorphs, which cannot all be
derived by mutation of a final consonant: \ird{-na/-ne},
\ird{-ime/-imně}, \ird{-é/-tě},
\ird{-ká/-ce}, and \ird{-ku/-cuje}. Their extra consonants, vowel
replacements, or quantity relations belong to the individual construction.
Thus linked \ird{mužne} preserves /n/ within the ending \ird{-ne};
it differs from adverbial \ird{mužně}, where PAL \ird{-ě} selects /ɲ/.
Contemplative \ird{mužnáž} instead gives \ird{mužnáže} by regular
linker suffixation: PAL preserves the input postalveolar /ʒ/.
The difference between free final [ʃ] and prevocalic [ʒ] follows
ordinary final devoicing (§\,\ref{sec:lenition}).


The linker selects the front grade described in
§\,\ref{sec:suffix-classes}. It shares the coronal and labial PAL mappings
with front-coded nominal endings, while its class-assigned velar output /ts/
differs from the /tʃ/ selected by mixed-grade irrealis \ird{-ěv}.
Plain-grade conditional \ird{-ic} and restricted velar-front \ird{-ir}
have different specifications. These four constructional classes combine two
instructions: whether the coronal and labial PAL mapping applies, and which
velar output is selected. The linker combines PAL with /k/ $\rightarrow$ /ts/;
irrealis \ird{-ěv} combines PAL with /k/ $\rightarrow$ /tʃ/;
\ird{-ir} selects /k/ $\rightarrow$ /ts/ without PAL; and \ird{-ic}
selects neither mutation. Surface vowel quality does not determine either
instruction. At each licensed boundary the mapping takes the selected input
consonant once; it does not feed its own output back into the same operation.

More specifically, the grade targets the final consonant of the inflected
input at its own boundary. It does not reach through an intervening aspect or mood exponent
to mutate the thematic consonant of the bare root. Where the evidential
intervenes, that input ends in its /l/, so the linker selects /ʎ/, written
\ird{lě} before its vowel, and the inner aspect or mood ending remains intact
(§\,\ref{sec:vm-evid}).

For input /k/, the front grade selects /ts/, spelled \ird{c}.
In \ird{pjášček} $\rightarrow$ \ird{pjáščice}, perfective
\ird{-ik-} is selected first, and its final velar then mutates.
The resulting /ts/ does not re-enter PAL at the same boundary.

Input /t\,d\,n\,l/ selects [c\,ɟ\,ɲ\,ʎ], whose sequences with
linker /ɛ/ are written \ird{tě\,dě\,ně\,lě}; input /s\,z\,ts/ selects /ʃ\,ʒ\,tʃ/, written
\ird{š\,ž\,č}.
Postalveolars already present within the form are unaffected.

Labials /p\,b\,m\,v/ acquire secondary palatalisation, marked by
\ird{ě} before /ɛ/, as in \ird{pjaštímě} and \ird{pjáščevě}.
Retrospective \ird{pjáščav} accordingly gives \ird{pjáščavě},
with prevocalic [ʋʲ] at the linker boundary.
Rhotic /r/ remains a rhotic: \ird{pjaštir} gives \ird{pjaštire},
without obligatory secondary palatalisation. Optional phonetic [rʲ]
is not an additional PAL output (§\,\ref{sec:soft-coronals}).

Consonant mutation precedes coda vocalisation and prevocalic approximant
realisation. The linked soft lateral remains [ʎ], whereas prevocalic /v/ carries its secondary
palatal articulation on [ʋʲ]. Coda /l/ and /v/ vocalisation is determined
by the syllabification of the complete linked form, before final devoicing
applies to its remaining obstruents
(§§\,\ref{sec:l-vocalisation} and~\ref{sec:v-realisation}).

The mappings apply to the boundary input once. Listed ending allomorphs in
Table~\ref{tab:linker-paradigm} supplement these productive correspondences;
they do not extend the rule to every segment found inside an ending.

The choice of the finite ending precedes mutation. In directly linked
perfectives without intervening Evid, the linker selects
bound \ird{-ik-}, after which \ird{k} becomes
\ird{c}: \ird{pjášček} $\rightarrow$ \ird{pjáščik-e}
$\rightarrow$ \ird{pjáščice}. The lexical root is unchanged. Transitive
\ird{pjaštník} already has the high-vowel exponent and gives
\ird{pjaštníce}. This suffixal correspondence differs from root raising in
\ird{zděk-}/\ird{zďíčev}. The perfective suffixal height relation
preserves the class-selected quantity, including long \ird{-ék}/\ird{-ík-};
irrealis G3 changes both height and quantity at its exposed root target.
Neither operation affects every \ird{e/ě} before a front-vocalic suffix.

Other linked aspect and mood forms in Table~\ref{tab:linker-paradigm} contain
construction-specific ending allomorphs. Final-vowel replacement or quantity
change in such pairs as \ird{mužna}/\ird{mužne} and
\ird{pjašté}/\ird{pjaště} belongs to those listed exponents, not to a
productive shortening or raising of the root.

Beyond these productive correspondences, conjunctive forms whose ending
correspondence exceeds ordinary mutation are listed with their aspect or mood paradigm. Thus final /ɛ/ in
\ird{pjaště}, written \ird{ě} after /c/, corresponds to free prohibitive \ird{-é}, while
\ird{pjaštímně} contains long \ird{-ímně}, corresponding to
Class~B \ird{-íme}; root-strong or lexical-long classes retain the
short \ird{-ime}/\ird{-imně} pair. These
forms require construction-specific allomorphs, even though they occupy the same
linker position as productive \ird{-ice}. Their selection requires the
corresponding finite ending at the linker boundary: \ird{pjáščekal}
instead yields \ird{pjáščekalě}, and \ird{pjaštímel} yields
\ird{pjaštímelě}. The evidential intervenes between Asp and the outer
linker, which therefore does not select inner \ird{-ik-} or
\ird{-imně}.

In prosodic terms, the linker \ird{-e} belongs to the same minimal lexical
prosodic word as its inflected host. Once the selected stem grade and the resulting syllables
have been established, primary stress falls on the first syllable of that
domain (§\,\ref{sec:stress}). Outer suffixation consequently does not move
primary stress to the linker. Subsidiary footing is assigned over the complete
surface word, including the linker, so that suffixation may supply the final
syllable needed to license a preceding nonfinal subsidiary head
(§\,\ref{sec:secondary-stress}). An integrated prefix remains inside the
domain and bears initial stress where present; an external negative
\ird{za-} does not enter it. Stress neutrality under suffixation therefore
does not imply primary stress fixed to the bare lexical root.

\subsection{Distribution and clause linking}

Within the verbal complex, \ird{-e} attaches outside all finite inflectional
morphology. The full template is:

\begin{center}
\Stem{} + \synhead{Asp.Align}.\textsc{align} + (\textsc{mood}) +
(\textsc{evid}) + \ird{-e}
\end{center}

\noindent The linker is therefore exterior to \synhead{EvidP}$_\text{inner}$,
the outermost projection in the finite inflectional spine
(§\,\ref{sec:vm-evid}). This position is significant: if \ird{-e} were an
agreement or mood morpheme, it would be expected inside the inflectional spine.
Its consistent extra-clausal position supports an analysis in which it marks the
edge of the finite complex as it interfaces with a higher projection --- a
\synhead{CP}-boundary marker, not an inflectional head.

In clause linking, the linker appears in exactly three environments:

\begin{enumerate}
  \item \textit{The finite \ird{po}-complement} (§\,\ref{sec:vm-po-comp}): the
    linker marks the right edge of the embedded finite \synhead{TP} before the
    complement C head \ird{po} or \ird{poby}. It is obligatory here when the
    embedded predicate is verbal; an embedded finite verbal clause without the
    linker is ungrammatical before \ird{po}/\ird{poby}. With zero-copula
    predicates, where no finite verb is present to host \ird{-e}, the support
    form \ird{ce} appears before the complementizer instead. The two
    complementizers attach not to a bare proposition versus a bare
    \textit{whether}-clause, but to assertive versus non-assertive finite
    complements respectively; the non-assertive type includes embedded
    wh-content clauses as well as embedded polar issues (see §\,\ref{sec:poby}
    of Chapter~\ref{ch:embedding}).

  \item \textit{Finite conjunctive linkage}
    (§\,\ref{sec:finite-conjunctive-linkage}): the linker marks a dependent
    finite predication within a clause complex. The linked clause retains its
    own finite inflection, while its relation to the following clause is
    interpreted through the predicates and discourse context.

  \item \textit{The matrixless \ird{po}-construction}
    (§\,\ref{sec:complementation}): even when the matrix predicate is absent and
    \ird{po} functions as a discourse-level attribution marker, the linker
    \ird{-e} is retained on the finite verb. This shows that \ird{-e} is not
    licensed by the presence of a lexical matrix predicate but by the presence
    of the C-domain element \ird{po} itself.
\end{enumerate}

These licensing environments contrast with two contexts in which the linker
is absent. First, independent finite clauses lack \ird{-e}: when no C head embeds the \synhead{TP}, the linker
is not licensed. Second, converb forms --- the non-finite clause-linking
morphology of simultaneous, sequential, conditional, and other converbs --- do
not take \ird{-e}. This follows naturally: converbs are already non-finite (they
lack Asp.Align and Mood morphology), so they do not participate in the domain of
the finite-clause boundary that \ird{-e} marks.

\subsection{Syntactic analysis}
\label{sec:vm-tension-linker}

The linker \ird{-e} appears at the right edge of embedded finite clauses before
\ird{po}/\ird{poby} and also in the conjunctive clause-chaining function
described in the preceding subsection. The analysis adopted here assigns these
uses a common clause-edge function, supported by their shared suffix position and
consonant-mutation class. The two environments remain distinguishable
in the syntactic contexts that license the suffix.

The grammar adopts a unified \synhead{CP}-edge analysis.
Its core claim is that \ird{-e} spells out the feature \textsc{[+fin, +sub]} ---
roughly, `finite clause under subordination' --- on the embedded \synhead{TP}'s
outermost layer, and that this feature is valued whenever a C head (overt or
null) selects the finite \synhead{TP}. The front grade selected by
\ird{-e} is then part of the PF realisation of this morphosyntactic feature,
parallel to the corresponding nominal mutation class. The exponent's vowel
does not itself derive this specification.

The position of \ird{-e} outside \synhead{EvidP}$_\text{inner}$ in the template
supports this analysis directly: the linker is the last element added to the
verbal complex before the clause exits the finite inflectional domain, which is
exactly where a \synhead{TP}-edge or C-interface marker would be expected.

Morphological locality provides a further point of correspondence: the
evidential supplies a consonantal boundary between the inner finite ending
and the linker. The linked sequences \ird{-alě} and
postvocalic \ird{-lě} follow from PAL applied to evidential /l/,
without mutation or linker-conditioned allomorphy across Evid
(§\,\ref{sec:vm-evid}). This local realisation is compatible with the clause-edge analysis.

\noindent The shared mutation class and ending selection associate
\ird{-e} with the morphological system of the verb, even though the linker
occupies the clause edge outside finite inflection. These properties support its
analysis as a suffix.


\subsection{Finite conjunctive linkage}
\label{sec:finite-conjunctive-linkage}

A non-final finite clause in \ird{-e} retains its own aspect, mood, and
aspect-conditioned case frame, while depending on the following clause for
completion of the conjunctive construction. In the clause-edge analysis
adopted here, a null linking C selects the finite constituent and licenses
its \textsc{[+fin, +sub]} edge. This construction differs from
\ird{po}/\ird{poby} complementation in the licensing C and the relation to the
following clause: the linked predication contributes to a clause complex
rather than serving as the selected content of a reporting, attitude, or
interrogative predicate.

The following constructed pair preserves the progressive following clause
while changing the linked predicate's aspect and alignment:
\pex\label{ex:finite-linking-alignment}%
\a
\begingl
\gla Marek lúbcí pjaštímně Janek znóvime.//
\glb Marek.\Dir{} apple.\Ind{} eat-\Prog{}-\Lnk{} Janek.\Dir{} study-\Prog{}//
\glft \trsl{Marek is eating the apple, and Janek is studying.}//
\endgl
\a
\begingl
\gla Lobek Marku pjaštníce Janek znóvime.//
\glb apple.\Dir{} Marek.\Trs{} eat-\Trs{}-\Pf{}-\Lnk{} Janek.\Dir{} study-\Prog{}//
\glft \trsl{Marek ate the apple, and Janek is studying.}
\lingcontext{the completed meal precedes the studying in the stated context}//
\endgl
\xe
The linked progressive in (a) has a direct actor and indirect patient; the
linked transitive perfective in (b) has a direct patient and transitive-case
agent. Neither copies the following clause's progressive morphology or case
frame. Finite forms select the same linker allomorphs as in complementation:
\ird{pjaštíme} gives \ird{pjaštímně}, while \ird{pjaštník} gives
\ird{pjaštníce}. By contrast, a converb licenses its own nominative-domain
frame without retaining finite viewpoint (§\,\ref{sec:converb-chains}).

The linked clause's finite aspect is interpreted locally within the complex,
with temporal anchoring supplied by the discourse; it does not require an
independent tense contrast absent from the verbal paradigm. Sequential,
circumstantial, or consequential relations may follow from the predicates and
context, although \ird{-e} itself neither specifies bounded anteriority like
\ird{-u} nor asserts a causal relation. The sequence in (b) illustrates a
chronological context rather than a separate lexical meaning of the linker.

Inner evidential morphology likewise remains inside the linked finite clause.
In established \ird{pjáščekalě}, Evid intervenes between the perfective and
linker, whose front grade targets evidential \ird{l}; the resulting form does
not replace the inner perfective with directly linked \ird{-ice}
(§\,\ref{sec:vm-evid}). This morphological independence permits the linked
predication to bear its own inner evidential specification. It does not by
itself fix the evidential's perspective holder or extend its scope over the
following clause. Outer stance, independently scoped questions or directives,
and separate illocutionary force require licensing beyond the suffix position;
their distribution within finite conjunctive complexes remains unspecified.

Finiteness consequently does not make the linked clause an autonomous
utterance. An affirmative complex can commit the speaker to both described
situations, but \ird{-e} neither supplies independent assertion nor assigns
a complete intonational contour to its host. Phrasing depends on the
integration and information structure of the complex, whereas an independent
sentence introduced by \ird{nu} can follow a full utterance boundary
(§\,\ref{sec:spoken-nu}).

Akkadian \textit{-ma} offers a comparison in which weakly specified finite
linkage supports several contextual relations under constraints on temporal
ordering \cite[58--60]{deutscher2009}. Its linked clauses are main clauses,
whereas Iridian \ird{-e} retains the dependent clause-edge analysis specified
here. The semantic comparison therefore does not identify their syntactic
structures or transfer Akkadian chronological restrictions to Iridian.


\section{Clausal complementation}\label{sec:complementation}

\subsection{The propositional/event-complement distinction}
\label{sec:th-compl-distinction}

Iridian distinguishes two major types of clausal complementation according to
the semantic type of the embedded content: \term{propositional complementation}
and \term{event complementation}. This distinction, which parallels the
broad cross-linguistic division between that-clause and to-infinitive
complements in English, is encoded formally by different constructions in
Iridian.

Propositional complementation is expressed by the finite \ird{po}-clause:
the embedded clause remains fully finite, retaining all verbal categories
of an independent clause, and is subordinated by the linker \ird{-e}
followed by \ird{po}. Event complementation is expressed by the purposive
converb: the embedded verb is non-finite and the construction is semantically
of the controlled or prospective type.

\begin{description}[leftmargin=2em, labelwidth=!, labelindent=0em]
\item[\upshape Propositional complement (\ird{po}):] used with predicates of
	saying, reporting, thinking, knowing, believing, remembering, inferring,
	perceiving a proposition, and evaluating a proposition. The complement
	presents a state of affairs as held to be true (or false, or probable,
	etc.).\ by the matrix subject.
\item[\upshape Event complement (purposive converb):] used with predicates of
	wanting, trying, intending, beginning, going (purposively), and related
	predicates of desire, attempt, inception, and motion. The complement
	presents an event as projected, volitional, incipient, or
	purpose-oriented rather than as a proposition held to be true.
\end{description}

\noindent The distinction is not merely semantic: it has formal consequences.
Propositional complements are always finite; event complements are always
non-finite. In languages where the distinction collapses (as it does for many
predicates in colloquial registers), Iridian preserves it morphologically.

\subsection{The finite \ird{po}-construction}

The ordinary means of embedding a full proposition is the finite
\ird{po}-clause. Its form is:

\begin{center}
\term{[finite embedded clause + linker \ird{-e}] \ird{po}}
\end{center}

The embedded clause remains fully finite. It preserves the same verbal
categories as an independent clause, including aspect, the alignment
distinctions associated with aspect, voice, and, where required by the sense,
mood. The construction is therefore not a reduced verbal form, but a true
finite clause made subordinate by the linker \ird{-e} and the complement
marker \ird{po}.

If the embedded predicate is a zero-copula predication and thus lacks an overt
finite verbal host, Iridian inserts the support form \ird{ce} immediately
before \ird{po}. This \ird{ce} is not a general copula but a clause-support
form used only to license the complement boundary where no finite verb is
available to bear the ordinary linker. Thus verbal predicates take
\ird{-e}~+\ird{po}, while zero-copula predicates take \ird{ce}~+\ird{po}.

An independent clause such as the following:

\ex
\begingl
\gla Lobek Marku pjaštník.//
\glb apple.\Dir{} Marek.\Trs{} eat-\textsc{trs}-\Pf{}//
\glft \trsl{Marek ate the apple.}//
\endgl
\xe

may be embedded without loss of finiteness:

\ex
\begingl
\gla Lobek Marku pjaštníce po Janku dálnik.//
\glb apple.\Dir{} Marek.\Trs{} eat-\textsc{trs}-\Pf{}-\Lnk{} \Compl{} Janek.\Trs{} say-\textsc{trs}-\Pf{}//
\glft \trsl{Janek said that Marek ate the apple.}//
\endgl
\xe

\ex
\begingl
\gla Marek doktor ce po Janku dálnik.//
\glb Marek.\Dir{} doctor.\Dir{} \textsc{supp} \Compl{} Janek.\Trs{} say-\textsc{trs}-\Pf{}//
\glft \trsl{Janek said that Marek is a doctor.}//
\endgl
\xe

The same construction is used after predicates of saying, reporting,
thinking, knowing, believing, remembering, inferring, perceiving a
proposition, and evaluating a proposition. In all such cases the complement
is construed as a full proposition, not merely as a projected event.

The \ird{po}-clause may also occur without an overt matrix predicate, while
retaining the finite complement structure. In these matrixless uses, the
proposition is presented as attributed content rather than asserted directly by
the speaker. According to context, this may
yield a reportive, hearsay, echoic, reflective, distancing, ironic, or
reservation-marking effect.

\ex
\begingl
\gla Lobek Marku pjaštníce po.//
\glb apple.\Dir{} Marek.\Trs{} eat-\textsc{trs}-\Pf{}-\Lnk{} \Compl{}//
\glft \trsl{Marek ate the apple, I hear.}//
\endgl
\xe

The precise omitted frame need not be supplied in every instance. One may
understand such clauses loosely as equivalent to \textit{it is said}, \textit{I
hear}, \textit{one gathers}, \textit{so they say}, or the like; but
grammatically the essential fact is that the proposition is cast as attributed
rather than directly asserted.

\subsection{The po-clause as a nominal \synhead{CP} complement}
\label{sec:vm-po-comp}

The finite propositional complement (\ird{po}-clause) is a \synhead{CP} that
occupies the complement position of the matrix predicate. The prototypical
structure of the complement \synhead{CP} is:

\begin{center}
\term{[\synhead{CP} [\synhead{TP} embedded-clause-\ird{e}] \ird{po}]}
\end{center}

\noindent \ird{Po} is the head of the complement \synhead{CP} (a
complementizer), and the embedded finite clause is the \synhead{TP} whose
head is marked by the linker \ird{-e}. In head-final structure,
the C head follows its complement \synhead{TP}:

\begin{center}
\begin{forest}
  syntax tree
  [\synhead{CP}
    [{\synhead{TP}\\\textit{embedded clause}\\\textit{(verb + -e)}}]
    [{C\\\ird{po}}]
  ]
\end{forest}
\end{center}

\noindent This \synhead{CP} is then the complement of the matrix verb:

\begin{center}
\begin{forest}
  syntax tree
  [\synhead{VP}
    [{\synhead{CP}\\{[\textit{emb-clause}-\ird{e}] [\ird{po}]}}]
    [V]
  ]
\end{forest}
\end{center}

\noindent Head-finality places the \synhead{CP} complement to the left of the
matrix V. The construction-specific placement of the whole selected clause
before ordinary matrix arguments yields the following surface order:

\begin{center}
\term{[embedded-\synhead{TP} + \ird{-e}] [\ird{po}] [matrix arguments] [matrix V]}
\end{center}

\noindent This is the order illustrated (§\,\ref{sec:complementation}).
The complement \synhead{CP} originates in [Complement, matrix-\synhead{VP}],
where head-finality places it before V. The selected finite CP moves as a whole to the construction-specific
external edge of \synhead{OrdP}, before the pivot and residual ordinary
constituents but below marked information structure. This placement is licensed
by a separate requirement for selected finite complements on Ord, applied
after ordinary placement while OrdP remains open. Its internal finite structure
is completed independently, and its internal arguments never become goals of
the matrix placement search. The lower CP occurrence preserves selection by V.

When the embedded predicate is a zero-copula predication, however, no overt
finite verb is available to host the boundary suffix \ird{-e}. In that special
case Iridian inserts the support form \ird{ce} immediately before \ird{po}
(and, in the corresponding non-assertive construction, before \ird{poby}).
This \ird{ce} is best analysed not as a general copular verb but as a
clause-support exponent used to realize the complement edge for otherwise
verbless finite predications. The resulting pattern is therefore:

\begin{center}
\term{[\synhead{CP} [small clause + \ird{ce}] \ird{po}]}
\end{center}

A recoverable matrix argument may remain unexpressed when the discourse
 identifies its referent. The propositional complement and an overt matrix
 topic occupy distinct positions, so the presence or weight of the
 \synhead{CP} does not syntactically exclude an overt topic. Null realisation
 follows the general recoverability conditions
 (§\,\ref{sec:tc-topic-drop}), independently of complement placement.


A finite complement and its matrix may form one intonational phrase, within
 which a weaker phonological boundary can mark the end of the complement
 (§\,\ref{sec:prosodic-domains}). Continuation pitch or modest final
 lengthening can make that edge audible without a pause; greater weight or
 rhetorical separation permits an intonational boundary instead. Neither the
 linker nor the complementiser requires an independent nuclear contour, so
 their clause-edge function does not entail a matching major prosodic edge
 (§\,\ref{sec:boundary-phenomena}).

The analysis of \ird{po} as a C head (and not, for example, as an
adposition or particle) is supported by the following diagnostics:
\begin{enumerate}
\item The embedded clause is fully finite (retains aspect,
  alignment and mood morphology); this is what is
  expected of a \synhead{CP} complement, not a reduced \synhead{VP}.
\item \ird{Po} appears after the embedded clause-final verb
  (the right edge of \synhead{TP} in a head-final \synhead{CP}), in exactly the
  position predicted for a head-final C.
\item \ird{Poby} (the non-assertive counterpart) parallels \ird{po}
  structurally and has the semantics of a complementizer: it
  types the embedded clause as an indirect interrogative
  (\textit{whether} or embedded wh-question) rather than an
  assertive declarative. Complementizer alternation is a
  \synhead{CP}-level phenomenon.
\item The linker \ird{-e} on the embedded verb is the edge-marking
  morpheme at the \synhead{TP}-\synhead{CP} boundary; where no verbal host is
  present because the predicate is zero-copular, the support form \ird{ce}
  appears instead. This is consistent with analysing the complement edge as a
  suffixal boundary marker on finite verbal clauses and as a support position
  on finite verbless predications (§\,\ref{sec:vm-tension-linker}).
\end{enumerate}

\subsection{Negative concord in finite complements}
\label{sec:complex-negative-concord}

Each finite clause forms its own negative-concord domain. A negative
indefinite inside a finite \ird{po}- or \ird{poby}-complement requires a
negative operator and the appropriate negative realisation within that
embedded clause; matrix negation cannot license an embedded indefinite
across the finite boundary. Conversely, a negative embedded clause does not
license a negative indefinite in an affirmative matrix clause. Within either
domain, several concord associates share one operator and express one
semantic negation (§\,\ref{sec:vm-neg}); a negative operator in each of two
clauses introduces two semantic negations, whose scope depends on the
embedding predicate.

\pex\label{ex:embedded-negative-concord}
\a
\begingl
\gla Deno Marku zapjaštníce po Janku dálnik.//
\glb nothing.\Dir{} Marek.\Trs{} \Neg{}-eat-\textsc{trs}-\Pf{}-\Lnk{} \Compl{} Janek.\Trs{} say-\textsc{trs}-\Pf{}//
\glft \trsl{Janek said that Marek ate nothing.}
  \lingcontext{one negative operator, inside the finite complement}//
\endgl
\a
\begingl
\gla *Deno Marku pjaštníce po Janku zadálnik.//
\glb nothing.\Dir{} Marek.\Trs{} eat-\textsc{trs}-\Pf{}-\Lnk{} \Compl{} Janek.\Trs{} \Neg{}-say-\textsc{trs}-\Pf{}//
\glft \trsl{Intended: Janek did not say that Marek ate anything.}
  \lingcontext{matrix negation cannot license embedded \ird{deno}}//
\endgl
\xe

In (\ref{ex:embedded-negative-concord}a), the direct-case patient \ird{deno}
is licensed by embedded negation, while \ird{Marku} is the transitive-case
agent of the perfective predicate. Linked \ird{pjaštníce} follows the
established perfective linker paradigm (§\,\ref{sec:vm-tension-linker}). The contrast
concerns concord licensing rather than the choice of complementizer:
non-assertive \ird{poby}, interrogative \ird{by}, and an affirmative
irrealis do not themselves provide a negative operator.

\subsection{Prospective and controlled event complements}

The finite \ird{po}-construction must be kept distinct from the complement used
for intended, attempted, controlled, or projected events. This second
construction belongs to the domain of the purposive converb. It is employed with
predicates such as \textit{want}, \textit{try}, \textit{intend}, \textit{go in
order to}, \textit{begin}, and similar verbs whose complement denotes not a
proposition as such, but an event regarded as prospective, volitional,
incipient, or purposive.

In this use the complement is non-finite. It does not present an independent
state of affairs as true, reported, believed, or remembered, but rather an event
toward which the matrix predicate is directed. The distinction is therefore one
between propositional complementation and event complementation. The former is
the province of finite \ird{po}-clauses; the latter, of the purposive converb.

\subsection{The purposive converb complement}
\label{sec:vm-purp-comp}

The purposive converb (\ird{-it}) serves as the complement of
desiderative, volitional, modal and motion predicates. Unlike the
\ird{po}-clause, it lacks finite viewpoint and mood, although its internal
arguments receive the NOM-domain cases licensed by the converb construction.
Its lower predicate can contain applicative and causative morphology; the
purposive head Conv, realised by \ird{-it}, closes this constituent without
introducing finite Asp.Align (§\,\ref{sec:vm-converb-aspless}).

The purposive complement is analysed as \synhead{ConvP}, whose head selects
the lower \synhead{vP}/\synhead{VoiceP}, rather than as a finite
\synhead{CP}. It denotes a prospective event description, while a
\ird{po}-complement presents a proposition. In head-final structure, Conv
follows its verbal complement and the entire purposive constituent precedes
the matrix predicate:

\begin{center}
\term{[matrix pivot] [\synhead{ConvP} [lower verbal predicate] \ird{-it}] [matrix V]}
\end{center}

\noindent The implicit participant of the purposive complement normally
shares the reference of the matrix absolutive pivot
(§\,\ref{sec:switch-reference}). In the ordinary intransitive-motion
examples, that pivot is the mover, so the dependency yields a same-subject
reading (§\,\ref{sec:purp-motion}). This reference relation is independent
of aboutness topichood and does not require an additional phonologically null
finite Asp head.
\section{The non-assertive complement \ird{poby}}\label{sec:poby}

\subsection{The \ird{po} / \ird{poby} opposition}

The fundamental contrast between \ird{po} and \ird{poby} is one of assertive
versus non-assertive complementation. \ird{Po} is the assertive complementizer:
it introduces a finite embedded clause presented as propositional content ---
something said, believed, remembered, known, or otherwise treated as an
assertable proposition. \ird{Poby} is the non-assertive complementizer: it
introduces a finite embedded clause presented as an issue or indirect question,
not as directly asserted content. Synchronically it is best treated as a
single lexicalized complementizer, not as a compositional sequence of
\ird{po}~+~\ird{by}; the independent interrogative particle \ird{by} remains the
overt exponent of Pol\textsuperscript{0} inside the embedded clause where that
head is overtly realized. The form in both cases is:

\begin{center}
\term{[finite embedded clause + linker \ird{-e}] \ird{po} / \ird{poby}}
\end{center}

The embedded clause remains fully finite in both constructions. Both \ird{po}
and \ird{poby} are therefore best analysed as heads of \synhead{CP} selecting a
clause that may itself project \synhead{PolP}; the complementizer contrast is a
\synhead{CP}-level phenomenon (§\,\ref{sec:vm-tension-linker} in
this chapter).

Neither non-assertive \ird{poby} nor clause-internal interrogative \ird{by}
licenses the negative-indefinite series without local negation
(§\,\ref{sec:complex-negative-concord}).

With zero-copula predicates the same contrast is preserved, but the support form
\ird{ce} appears instead of verbal \ird{-e}: \term{[small clause + \ird{ce}]
\ird{po}/\ird{poby}}.

\subsection{Two complement types under \ird{poby}}

The non-assertive framing of \ird{poby} covers two core complement types. These
are not unrelated constructions but two values of the same broader indirect
interrogative category.

\noindent\textit{Type~1: Embedded polar issue (\textit{whether}-complement).}\quad  The
first and most central use embeds an unresolved polar question: whether a
proposition holds. This is the closest Iridian equivalent to English
\textit{whether}. The embedded clause is a full finite clause without a
wh-operator. After verbs of asking, wondering, doubting, or deciding, \ird{poby}
marks the embedded content as an open issue rather than an asserted fact.

\pex
\a
\begingl
\gla Lobek Marku pjaštníce poby Janku houstnik.//
\glb apple.\Dir{} Marek.\Trs{} eat-\textsc{trs}-\Pf{}-\Lnk{} \Compl{} Janek.\Trs{} ask-\textsc{trs}-\Pf{}//
\glft \trsl{Janek asked whether Marek ate the apple.}//
\endgl
\a
\begingl
\gla Marek ščénice poby Janku houstnik.//
\glb Marek.\Dir{} arrive-\Pf{}-\Lnk{} \Compl{} Janek.\Trs{} ask-\textsc{trs}-\Pf{}//
\glft \trsl{Janek asked whether Marek arrived.}//
\endgl
\xe

\noindent\textit{Type~2: Embedded wh-content question.}\quad
The second type embeds a content question: who came, what he bought, where they
went. Here the embedded clause contains a wh-operator from the inventory in
Table~\ref{tab:nm-interrogatives} (Chapter~\ref{ch:nouns}), and \ird{poby} marks
the clause as non-assertively embedded. The embedded wh-phrase participates in ordinary placement without dedicated
wh-fronting in the neutral construction; optional fronting within the embedded
clause's own left periphery signals marked focus, consistent with the general
principle established in §\,\ref{sec:th-wh-interrogation}.

\pex
\a
\begingl
\gla Jede ščénice poby Janku houstnik.//
\glb who.\Dir{} arrive-\Pf{}-\Lnk{} \Compl{} Janek.\Trs{} ask-\textsc{trs}-\Pf{}//
\glft \trsl{Janek asked who arrived.} \lingcontext{without dedicated wh-fronting}//
\endgl
\a
\begingl
\gla Marek jena mulše poby Janku houstnik.//
\glb Marek.\Dir{} where live-\Cont{}-\Lnk{} \Compl{} Janek.\Trs{} ask-\textsc{trs}-\Pf{}//
\glft \trsl{Janek asked where Marek lives.} \lingcontext{without dedicated wh-fronting}//
\endgl
\xe

\noindent\textit{Direct quoted questions.}\quad
Directly quoted questions are not ordinarily introduced by \ird{poby}. A direct
quote preserves the quoted clause's own illocutionary force and is therefore
framed by the ordinary quotative complementizer \ird{po}; if the quoted clause
is a polar question, the interrogative particle \ird{by} remains inside the
quoted clause itself. The function of \ird{poby} is instead to embed the
corresponding question indirectly as an issue-clause. The contrast is therefore
between direct quotation under \ird{po} and indirect question embedding under
\ird{poby}, not between two alternative quotative particles.

Where the directly quoted clause is zero-copular, \ird{ce} supplies the support
position before \ird{po}. Thus a quoted predicational clause such as
\ird{Marek doktor ce po} is structurally parallel to verbal
\ird{Marek ščénice po}: in the former \ird{ce} supports the complement edge,
while in the latter the finite verb bears \ird{-e}.

\subsection{Internal syntax of \ird{poby}-clauses}

In \ird{poby}-complements, the embedded clause projects \synhead{PolP}, but
neutral embedded content ordinarily uses null Pol\textsuperscript{0}. The
complementizer \ird{poby} already marks the clause as non-assertive and
issue-like; an embedded plain polar issue therefore does not require an
additional overt \ird{by} internally. Overt internal \ird{by} is reserved for
marked, echoic, or explicitly interrogative embedded readings, where the speaker
wishes to foreground the interrogative status of the embedded content rather
than presenting it as a background issue.

The architecture for the two core types is thus as follows. An embedded polar
issue has a fully finite embedded clause with null Pol\textsuperscript{0}: the
non-assertive status is contributed entirely by \ird{poby}. An embedded
wh-content clause contains the wh-operator plus null Pol\textsuperscript{0}. A
marked embedded echo question may contain overt internal \ird{by} alongside the
wh-operator, in a double-marking configuration that asserts the explicitly
interrogative, insisted, or surprised character of the embedding.

A directly quoted question is structurally different. There the
embedded clause is quoted under \ird{po}; interrogative force is carried by the
quoted clause's own Pol\textsuperscript{0} (overt \ird{by} or null, as in
matrix questions), and \ird{poby} does not appear.

This construction is found after predicates such as \textit{ask},
\textit{wonder}, \textit{inquire}, \textit{investigate}, \textit{check},
\textit{doubt}, \textit{decide}, and \textit{know whether}, together with other
predicates that take indirect interrogative complements. Predicates of the
\textit{say that}, \textit{believe that}, \textit{remember that}, or
\textit{know that} type require the assertive \ird{po}-construction; directly
quoted questions likewise pattern with \ird{po}, not with \ird{poby}.

\subsection{Indirect evidentiality and attributed content}
\label{sec:complex-evidential}

Finite clauses in Iridian are, in themselves, evidentially unmarked. The
ordinary declarative does not grammatically specify whether the speaker knows
the matter directly or indirectly, though in practice it is often compatible
with direct knowledge. Where the source of knowledge is explicitly indirect,
Iridian employs a special indirect evidential. This marks that the proposition
rests on report, inference, visible result, or post hoc discovery rather than on
firsthand evidence. In suitable contexts a mirative colour may arise from this
meaning, especially where the information is sudden or contrary to expectation;
but the grammatical value remains that of indirect evidence.

This category must not be confused with the matrixless \ird{po}-construction.
The indirect evidential concerns the source of knowledge. A matrixless
\ird{po}-clause concerns the mode of presentation in discourse. The two
frequently approach one another in use, since hearsay lies near the border of
both; but they are not the same thing.

Morphologically the indirect evidential belongs to the finite verbal complex. It
stands before the complement linker \ird{-e}, occupying the outermost position
in the finite inflectional spine. The order is therefore as follows:

\begin{center}
\textsc{stem} + \textsc{asp.align} + (\textsc{mood}) + (\textsc{ind.evid}) +
\ird{-e} + \ird{po}/\ird{poby}
\end{center}

The default evidential exponent is \ird{-al} after a consonant and
\ird{-l} after a vowel (§\,\ref{sec:vm-evid}). Both attach to the completed
finite form, preserving its stem grade, aspect and mood exponents, and vowel
quantity. When the clause takes the outer linker, its PAL instruction applies
to the evidential's final /l/, yielding \ird{-alě} or vowel-final base plus
\ird{-lě}. It does not reach through the evidential to select an inner linked
form such as \ird{pjaštníce} or \ird{pjaštímně}. Final evidential /l/
is realised as coda [w] in an independent form, whereas linked /ʎ/,
written \ird{lě} with the linker vowel, occupies the onset before /ɛ/; neither realisation changes the lexical
primary stress or vowel quantity of the host. Subsidiary footing is assigned
to the completed lexical word, including the evidential and any linker
(§\,\ref{sec:secondary-stress}).

The perfective \ird{pjaštník} consequently gives \ird{pjaštníkal} and
linked \ird{pjaštníkalě}, while progressive \ird{pjaštíme} gives
\ird{pjaštímel} and linked \ird{pjaštímelě}. The complete aspectual
paradigm is given in Table~\ref{tab:inner-evidential-paradigm}, and
Table~\ref{tab:evidential-combinations} includes Voice, causative, and mood
combinations. In each case the evidential intervenes between the completed
finite host and the linker, including when that host already contains several
valency-changing and inflectional exponents.

A finite evidential can stand independently; the linker appears only when the
clause's embedding or linking construction requires it:

\pex
\a
\begingl
\gla Lobek Marku pjaštníkal.//
\glb apple.\Dir{} Marek.\Trs{} eat-\Trs{}-\Pf{}-\textsc{ind.evid}//
\glft \trsl{Marek ate the apple, apparently.}//
\endgl
\a
\begingl
\gla Lobek Marku pjaštníkalě po Janku dálnik.//
\glb apple.\Dir{} Marek.\Trs{} eat-\Trs{}-\Pf{}-\textsc{ind.evid}-\Lnk{} \Compl{} Janek.\Trs{} say-\Trs{}-\Pf{}//
\glft \trsl{Janek said that Marek had apparently eaten the apple.}//
\endgl
\xe

When the speaker chiefly wishes to indicate that the information is second-hand,
inferred, or otherwise non-firsthand, the indirect evidential is the ordinary
means. When the speaker chiefly wishes to cast the proposition as attributed
discourse, with reportive, distancing, sceptical, ironic, or echoic effect, the
matrixless \ird{po}-construction is the ordinary means. The two may co-occur,
but such co-occurrence is marked, and yields a layered interpretation: a
proposition already marked as indirectly known is then further presented as
attributed content. This is suitable for rumour, mediated narration, gossip, and
similar discourse types, but should not be treated as the unmarked way of
expressing hearsay.

\subsection{Particle domains and attribution}
\label{sec:particle-domains-attribution}

The finite verb bearing \ird{-al}/\ird{-l} identifies the predication
marked for indirect evidence, whereas \ird{po} delimits attributed
content. In \ird{Lobek Marku pjaštníkalě po Janku dálnik}, the
evidential belongs to the embedded predication of eating; the matrix
predicate attributes that content to Janek. This distribution does not by
itself determine whether the indirect evidence is evaluated from Janek's
perspective or from that of the reporting speaker. The source of a report,
the holder of an inference, and the speaker presenting either are distinct
interpretive parameters, whose identification cannot be recovered from
suffixal order alone. No general rule of perspective shift for embedded
particles is specified here.

Tagalog reportative \textit{daw} illustrates a further distinction between
scope and projection. With modal \textit{baka} `maybe', Kierstead and
Martin distinguish a report that $p$ is possible, the possibility of a
report that $p$, and a projective reading in which a report that $p$
is maintained while $p$ itself remains merely possible. In their
contexts, \textit{daw} after the embedded predicate permits all three
readings, whereas its position immediately after \textit{baka} does not
permit either narrow reportative scope or that projective reading
\cite[§2, especially Table~1, pp.~330--332]{kiersteadmartin2012}.
The survival of reportative content outside the modal's effect therefore
does not establish reportative scope over the modal. This comparison
distinguishes interpretive possibilities without assigning the same
readings to Iridian \ird{či}, the inner evidential, or attributed
\ird{po}.

Outer \synhead{EvidP} and \synhead{SpeechActP} concern the clause
periphery, rather than the morphological boundary between the finite verb
and its linker (§\,\ref{sec:th-particle-framework}). Reportative
\ird{či} can co-occur with quotative framing
(§\,\ref{sec:class2-particles}), and the inner evidential can combine
with matrixless \ird{po} in the marked layered use described above.
These established combinations provide no general licence for several
outer particles, or for their occurrence in every finite complement.
Likewise, the retention of finite verbal categories under \ird{po} and
\ird{poby} does not establish the availability of every matrix
speech-act construction inside their complements. Matrix stance towards a
report, stance within attributed content, and stance towards the act of
asking an embedded question involve different domains; adjacency to the
complement does not choose among them. The unspecified combinations in
§\,\ref{sec:particle-combinations} retain that status across the
complement boundary.

\section{Relativisation}
\label{sec:th-relativization}

\subsection{Structure of the relative construction}
\label{sec:th-rel-structure}

A relative in \ird{-ový} precedes its overt head noun and identifies a
participant privileged by the embedded voice--aspect configuration
(§\,\ref{sec:nm-ovni}). The construction has no overt relative pronoun or
complementiser. Its relative head selects an aspectual predicate and projects
\synhead{RelP}, with a silent operator linked to the privileged participant
gap. In this analysis, \ird{-ový} is a distinct relative head, rather than a
converb or a case form of nominal \ird{-ou}.

\begin{center}
[\textit{indirect dependents + voice--aspect--Rel}]\quad[\textit{head noun}]
\end{center}

The perfective counterpart of \ird{pjáščikou} is
\ird{pjáščikový}, whose antipassive Voice selects the agent. Patient-oriented
\ird{-n-} gives \ird{pjaštníkový}; active continuous, progressive, and
contemplative relatives select agents without antipassive morphology.
Continuous Asp and Rel have the joint exponent \ird{-ový}, parallel to
continuous Asp and $n_{\mathrm{part}}$ in \ird{-ova}; progressive,
retrospective, and contemplative preserve their construction-selected bound
aspectual material before relative \ird{-ový}. Ordinary property attribution
uses the dedicated stative \ird{-a}/\ird{-á} construction. Continuous
stative \ird{-ový} is licensed only in an explicitly episodic, stage-level,
temporally delimited or contrastive context: \ird{mužnový nuž} can denote
a city beautiful during the relevant celebration, while \ird{múžna nuž}
is the ordinary property modifier (§\,\ref{sec:stative-attr}). The relative
operator abstracts over the state holder selected as the sole participant;
additional valency morphology must obey the same highest-frame selection as
other relatives. \synhead{StateP}, \synhead{BecomeP} and event-derived
\synhead{ResP}, where independently licensed, belong below Asp.
The selected constituent is actual
\synhead{AspP}, containing viewpoint/Asp.Align, orientation Voice, causative
and inner valency structure. It excludes the higher indicative Neg domain,
Phas, the Comm-field, Mood and clausal evidentiality, as in participant
nominals (§\,\ref{sec:nominalised}). In particular,
\ird{[lúbcí pjáščikový] pres} contains a positive eating description, while
an intended reduced modifier \trsl{the man who still/already/only/did not eat
the apple} cannot be obtained by inserting those clausal operators inside the
bracket. An operator outside the relative may modify the matrix predication
containing that man; its external scope must not be read as an operator inside
\synhead{AspP}. This selection restriction does not exclude negative finite
relatives: their finite clausal structure, where independently licensed, can
contain negation. The reduced \ird{-ový} construction supplies neither a
finite negative relative paradigm nor a universal restriction on relatives.

\subsection{Head-finality and the placement of relative clauses}
\label{sec:th-rel-headfinal}

The relative clause occupies a prenominal modifier position within the host
noun phrase, consistent with the order of attributives and demonstratives
elsewhere in the nominal system. No extraposition or rightward movement of
relative clauses occurs in standard Iridian.

The matrix clause assigns case to the head noun independently of its
interpretation inside the relative. Within the modifier, thematic dependents
receive indirect case through nominal relational licensing, just as they do
inside participant nominals. Matrix case therefore does not express the role
of the relative gap.

\subsection{The participant gap and accessibility}
\label{sec:th-rel-gap}

The silent operator binds the position privileged by embedded Voice and aspect;
the resulting participant property restricts the overt head noun. An agent
and a patient are consequently accessible through different voice--aspect
configurations rather than through unrestricted gapping of either finite
argument. Applicative morphology makes the promoted participant accessible by
the same selection mechanism. A causative above the applicative determines
selection from the derived frame: its imperfective active privileges the
causer, while its ergative configuration privileges a retained theme, then an
applied participant if there is no theme, or otherwise the causee. A lower
applied participant therefore becomes the gap only when the highest
configuration selects it (§\,\ref{sec:aspect-caus}). The overt head noun
receives its case from the matrix clause, independently of the case its
referent would bear in a finite counterpart of the relative.

The established construction does not specify a resumptive-pronoun strategy
or a productive means of relativising every postpositional or oblique relation.
Such positions cannot be added to the paradigm solely by assuming a general
cross-linguistic accessibility hierarchy; the described \ird{-ový} system
selects the privileged participant.

The restriction is construction-specific. Finite linked topicalisation
 permits a transitive-case agent while its patient remains direct, and finite
 content questions permit interrogative expressions with ordinary placement
 or marked wh-fronting
 (§§\,\ref{sec:tc-marked-topics}, \ref{sec:th-wh-interrogation}). A restriction
 on the participant abstracted by a reduced relative therefore supplies no
 language-wide ban on agent A$'$-dependencies. Table~\ref{tab:construction-accessibility}
 distinguishes these established constructions.

\begin{table}[htbp]
\footnotesize
\caption{Accessibility by construction. Selection within a reduced predicate
 is distinct from finite argument case and from high discourse dependencies.}
\label{tab:construction-accessibility}
\begin{tabularx}{\textwidth}{L L L}
\toprule
Construction & Established dependency & Structural limit \\
\midrule
Linked topic & Direct- or transitive-case argument linked to a lower position &
Case preserved; independent aboutness relation in \synhead{TopP} \\
Finite content question & Ordinary placement or marked wh-fronting &
Lower occurrence interpreted under Pol; marked focus lands above the ordinary
placement field \\
Reduced \ird{-ový} relative & Operator linked to the participant selected by
embedded Voice and aspect & \synhead{RelP} selects the reduced aspectual
predicate; no finite I-position required \\
Participant nominal \ird{-ou} & Abstraction over the participant selected by
embedded Voice and aspect & Nominal head and nominal case domain;
not finite wh-extraction \\
Hanging topic & External referent linked to resumption or a recoverable argument &
No case-preserving argument-movement dependency established \\
Cleft or coordination extraction & No general accessibility paradigm established &
No universal extraction rule inferred from the preceding rows \\
\bottomrule
\end{tabularx}
\end{table}

\noindent Relative \ird{-ový} and participant-nominal \ird{-ou} share
 participant-selection conditions but differ in their selecting heads and
 external distribution (§\,\ref{sec:nm-derived-participant}). Their
 indirect dependents belong to the reduced construction's nominal licensing
 domain, rather than reproducing the full finite case frame. Shared selection
 consequently does not establish a shared landing position, and neither
 construction requires movement into a finite \synhead{IP}. The overt head
 of a relative receives matrix case; a participant nominal instead forms a
 nominal argument in its own right.

\section{Adverbial particles and the clause periphery}
\label{sec:th-particle-integration}

The particle hierarchy described in §\,\ref{sec:adv-particles} constrains
the interpretation of complex clauses: particle ordering, scope, and fragment
licensing identify the left-peripheral positions within which interrogative and
evidential operators are interpreted.

Each finite clause constructs its own ordinary placement and particle fields,
so that the arguments and modifiers of an embedded finite predicate precede
its particles within the complement delimited by the linker and
\ird{po}/\ird{poby}. Movement of a selected clause as a whole preserves that
internal organisation; matrix placement does not extract arbitrary embedded
arguments across the finite boundary. Reduced relatives, participant nominals
and converbs retain their independently selected structures, without acquiring
the full finite placement field or additional Mood, Pol or stance projections.

The hierarchy is summarised in Table~\ref{tab:th-clause-spine} and
Figure~\ref{fig:th-left-periphery}. A wh-expression undergoes ordinary
placement before the particle field without dedicated wh-fronting, while its
lower occurrence supplies an interrogative variable in the interpretation
domain of Pol (§\,\ref{sec:th-multiple-wh}). Marked wh-fronting instead
places one expression in Spec,\synhead{FocP}, above the ordinary placement
field, accounting for the contrast in (\ref{ex:wh-four-combinations}) with
either overt or null Pol. Neither dependency requires an aboutness topic.
When all ordinary constituents are unpronounced, a particle may surface
initially; topicless presentations independently lack a high topic dependency
(§\,\ref{sec:tc-topicless}).

Locally head-initial \ird{by} precedes the lower comment-domain operators,
including \ird{jas}, although overt arguments precede both by ordinary
placement (§\,\ref{sec:by-pol-head}). Distinct syntactic positions do not
by themselves license every combination or determine the projection of
recurrent \ird{jas vy} in embedded issues. The marked combination of the
inner indirect evidential with matrixless \ird{po} likewise preserves the
distinction between indirect evidence and attributed presentation, whose
licensing does not follow from formal separation alone
(§\,\ref{sec:particle-domains-attribution}).

The related constructions are described in the following sections:

\begin{description}[leftmargin=2em, labelwidth=!, labelindent=0em]
\item[\upshape Left-periphery structure (§\,\ref{sec:th-left-periphery}):] see
	§\,\ref{sec:adv-particles} for the full particle inventory and
	§\,\ref{sec:particle-combinations} for combination restrictions and
	scope contrasts; the formal analysis is in
	§\,\ref{sec:th-particle-framework}, with Figures~\ref{fig:th-clause-fields}
	and~\ref{fig:th-left-periphery}.
\item[\upshape Polar questions (§\,\ref{sec:th-polar-questions}):] see
	§\,\ref{sec:by-pol-head} (Chapter~\ref{ch:clause-types}) for the
	structural analysis of \ird{by} as the overt exponent of
	\synhead{Pol\textsuperscript{0}}, its independence from the
	comment-domain operators, and the scope account of mirative and
	reportative polar questions.
\item[\upshape \textit{Wh}-interrogation (§\,\ref{sec:th-wh-interrogation}):] see
	§\,\ref{sec:th-wh-interrogation} for single- and multiple-wh patterns, the
	lower-occurrence licensing principle, and the typology of question types.
\item[\upshape Evidential particles (§\,\ref{sec:class3-particles}):] see
	§\,\ref{sec:class3-particles} (\ird{hladě}, \ird{kady}, \ird{oče})
	and §\,\ref{sec:class2-particles} (\ird{či}) and §\,\ref{sec:complementation} for the interface between
	evidential morphology and attributed discourse.
\item[\upshape Complementation (§\,\ref{sec:complementation}):] see
	§\,\ref{sec:complementation} for the full treatment of \ird{po}/\ird{poby},
	the purposive converb complement, and the matrixless \ird{po}-clause.
\item[\upshape Head-finality and subordination (§\,\ref{sec:th-subord-headfinal}):] see
	§\,\ref{sec:converb-chains} for the syntax and information structure of
	converb chains, and §\,\ref{sec:sequential}--§\,\ref{sec:purposive} for
	individual subordinate types.
\end{description}
